% ============================================
% Discussion
% ============================================
\section{Discussion}
\label{sec:discussion}


Our work demonstrates that effective interaction intensity, environmental feedbacks, and assembly-generated interaction structure govern community coalescence outcomes, determining both the dominant outcome type and the level at which selection operates. In both experiments and modeling, when effective interactions are strong, Dominance prevails and species persistence is correlated within parental communities, indicating community-level selection; when interactions are weak, Mixture predominates and species respond independently. Furthermore, predictability analyses suggest two regimes within Dominance. In the top-down regime, dominant taxa and environmental feedbacks, including pH modification, may bias winner identity. In the emergent regime, collective multi-species dynamics appear to govern outcomes, and dominant-species-based prediction is insufficient. These findings reconcile the Clements--Gleason debate: effective interaction intensity determines whether communities behave as cohesive units or loose species assemblages.

Our work provides the first experimental demonstration of alternative regimes in community coalescence, reconciling previously contrasting observations of community-level versus species-level selection. Prior studies have reported both outcomes: some found communities behaving as cohesive units\citep{Tikhonov2016, DiazColunga2022}, while others observed species responding independently\citep{VanderGucht2007, Vermeij1991}. Our results suggest that these conflicting observations may reflect differences in interaction strength across experimental systems. For instance, the absence of community-level selection in \textit{in vitro} gut microbiome coalescence experiments\citep{Goldman2025, Walton2025} is consistent with our framework. Rich media such as Brain Heart Infusion (BHI) are analogous to our Nutr$-$ condition, providing complex and diverse resources while lacking high concentrations of readily metabolizable carbon sources. This resource structure reduces both direct resource competition and pH-mediated interactions, thereby weakening interspecies interactions\citep{Ratzke2020} and favoring species-level rather than community-level dynamics. Likewise, interaction strength may vary systematically across ecosystems: microbial biofilms experience strong metabolic interactions due to high cell densities and shared resources\citep{Nadell2016}, whereas mobile macroscopic organisms interact more transiently across larger spatial scales. These differences may explain why community-level selection is more frequently observed in microbial coalescence studies. This framework provides a unified perspective: rather than asking whether communities are cohesive units, we should ask under what conditions they become so, with interaction strength emerging as a key determinant.


A growing body of work has identified interaction strength as a key control parameter governing diverse aspects of microbial community dynamics, including diversity \citep{Hu2022, Marsland2019}, stability \citep{Allesina2012, Ratzke2020}, coexistence \citep{Grilli2017}, and invasion outcomes \citep{Hu2025, Kurkjian2021}. Our results extend this claim to include coalescence, providing further evidence that interaction strength determines when communities exhibit collective behavior. Importantly, interaction strength is used here as a coarse-grained descriptor that can reflect the combined effects of multiple ecological processes, rather than as a claim about one specific mechanism. The parallel with priority effects is particularly instructive; recent work showed that assembly order shapes final composition only under strong interactions, whereas weak interactions lead to convergent outcomes regardless of assembly history \citep{Hu2025, Fukami2015}. Our coalescence results mirror this pattern: strong interactions allow parental-community identity to shape the outcome, while weak interactions yield convergent Mixtures. Together, these findings reinforce interaction strength as a useful predictor of history-dependent community dynamics across multiple ecological contexts.

Notably, and perhaps counterintuitively, our model shows that community-level selection-like patterns can emerge without explicitly including cooperative interactions between species, consistent with predictions from Tikhonov's resource-competition models \citep{Tikhonov2016} and random Lotka--Volterra theory \citep{May1972, Grilli2017}. Community-level selection is often attributed to cooperative mechanisms such as cross-feeding or division of labor. By contrast, the model suggests that assembly history and competitive interactions can be sufficient to generate correlated species fates during coalescence. Thus, cooperative and competitive mechanisms may represent distinct routes to community-level selection, and understanding how these routes interact remains an important direction for future work.

Accordingly, we use community-level selection as an outcome-level description of origin-correlated persistence during coalescence. This operational definition does not by itself identify the causal ecological or biochemical mechanisms, which may vary across systems. In nature, coalescence occurs in richer settings that may alter regimes and outcomes. Environmental heterogeneity (temperature, pH buffering, and other physicochemical factors) covaries with interaction strength and can reshape coalescence across habitats \citep{Rillig2015, Castledine2020, Rillig2017, Tropini2017}. In host-associated microbiomes, host filtering, priority effects, and spatial structure further influence successful colonization, making coalescence a useful framework for predicting and steering microbiome transfer \citep{Liu2024, Smillie2018, Rocca2021}. Our experiments focus on steady-state outcomes and thus do not capture temporal dynamics such as fluctuating resources, migration pulses, or host responses. Accounting for this broader ecological complexity will help build a more comprehensive picture of community-level selection across habitats and scales.
